% GENERATED FILE. Edit canonical lessons, then run npm run generate:books.
% canonical-source-hash: fnv1a64:dbe136407e185e85
% canonical-lessons: RU-C265-kak-proyti, RU-C265-idite-pryamo, RU-C265-review

\chapter{Asking the Way}
\label{ch:ru-asking-the-way}
\hlchaptermodality{265}

\begin{chapteropening}
\textbf{By the end of this chapter:} \emph{I can ask the way with \ru{Извините}, \ru{как} \ru{пройти} \ru{к} ...?, and follow an answer with \ru{прямо}, \ru{направо}, \ru{налево} and \ru{поверните}.}

Complete RU-C265-review: ask the way to two places in town and give one set of directions back.
\end{chapteropening}

\section[\texorpdfstring{\ru{как} \ru{пройти} \ru{к} …? (kak proytí k …?)}{как пройти к …? (kak proytí k …?)}]{\ru{как} \ru{пройти} \ru{к} … (kak proytí k …) — asking the way}

\label{lesson:RU-C265-kak-proyti}



\begin{quote}
Say \textquotedblleft{}excuse me\textquotedblright{}, \textquotedblleft{}where\textquotedblright{} and \textquotedblleft{}the station\textquotedblright{}. (\emph{Izvinítye. Gde. Vokzál.}) Then \textquotedblleft{}I must\textquotedblright{}: a man and a woman. (\emph{Ya dólzhen. Ya dolzhná.})

Asking \emph{\ru{где} \ru{вокзал}} works. Russians also ask for the way itself.
\end{quote}

\subsection*{You'll want to know: \ru{как} \ru{пройти}}
\begin{quote}
\textbf{\ru{Извините}, \ru{как} \ru{пройти} \ru{к} \ru{вокзалу}?} — \emph{izvinítye, kak proytí k vokzálu?} — Excuse me, how do I get to the station?
\end{quote}

\textbf{\ru{Пройти}} is the perfective of a going-verb: to get somewhere on foot, through the streets, and arrive. It stays in the infinitive, with no \textquotedblleft{}I\textquotedblright{} in the question: literally \textquotedblleft{}how to get through to the station\textquotedblright{}.

\begin{grammarlens}[title={Grammar Lens: \ru{к} and the \textquotedblleft{}to\textquotedblright{} form}]
\textbf{\ru{к}} — \emph{k} — \textquotedblleft{}towards, to\textquotedblright{}. The place after it takes the same \textquotedblleft{}to\textquotedblright{} form as \emph{\ru{маме}} and \emph{\ru{брату}}:

\begin{quote}
\textbf{\ru{к} \ru{вокзалу}}, \textbf{\ru{к} \ru{музею}} — \emph{k vokzálu, k muzéyu} — a masculine noun ending in a consonant takes \textbf{-\ru{у}}, and one in \emph{-\ru{й}} takes \textbf{-\ru{ю}}. \textbf{\ru{к} \ru{остановке}} — \emph{k ostanóvke} — most nouns in \emph{-\ru{а}} take \textbf{-\ru{е}}. \textbf{\ru{к} \ru{площади}} — \emph{k plóshchadi} — \emph{\ru{площадь}}, in \emph{-\ru{ь}}, takes \textbf{-\ru{и}}. \textbf{\ru{к} \ru{метро}} — \emph{k metró} — \emph{\ru{метро}} stays as it is.
\end{quote}
\end{grammarlens}

\subsection*{The exchange}
\begin{quote}
\textbf{— \ru{Извините}, \ru{как} \ru{пройти} \ru{к} \ru{музею}? \ru{Это} \ru{далеко}?} \textbf{— \ru{Нет}, \ru{это} \ru{близко}. \ru{Музей} \ru{рядом}.}
\end{quote}

\emph{Éto dalekó?} — \textquotedblleft{}Is it far?\textquotedblright{} \emph{Nyet, éto blízko. Muzéy ryádom} — \textquotedblleft{}No, it's near. The museum is close by.\textquotedblright{}

\begin{cousinweb}
\textbf{\ru{пройти}} is \emph{\ru{про}-}, \textquotedblleft{}through\textquotedblright{}, with \textbf{\ru{идти}}, \textquotedblleft{}to go\textquotedblright{}, which EtymDB 2.1 traces through Old East Slavic \emph{iti} and Proto-Slavic \emph{*jĭti} to Proto-Indo-European \emph{*h\textsubscript{1}éyti}, \textquotedblleft{}to go\textquotedblright{}, the root of Latin \emph{īre}. \textbf{\ru{к}} goes back to Proto-Slavic \emph{*kŭ}, which EtymDB links to a Proto-Indo-European \emph{*kom}, \textquotedblleft{}next to, with\textquotedblright{}.
\end{cousinweb}

\subsection*{Your turn}
\begin{itemize}
\raggedright
  \item \emph{Say it:} \emph{Izvinítye, kak proytí k vokzálu? Éto dalekó?}
  \item \emph{Say it:} the same to the museum, the stop and the square: \emph{k muzéyu}, \emph{k ostanóvke}, \emph{k plóshchadi}
  \item \emph{Answer:} \emph{Kak proytí k metró?} \emph{Éto blízko} or \emph{Éto dalekó}
\end{itemize}

\begin{tcolorbox}[breakable,colback=teal!4,colframe=teal!35!black,title={Before you move on}]
\textquotedblleft{}How do I get to the station?\textquotedblright{} (\emph{\textbf{\ru{Как} \ru{пройти} \ru{к} \ru{вокзалу}?}}) What does \emph{\ru{к}} do to \emph{\ru{остановка}}? (\emph{\textbf{\ru{К} \ru{остановке}.}}) And to \emph{\ru{метро}}? (\textbf{Nothing:} \emph{\textbf{\ru{к} \ru{метро}}}.)
\end{tcolorbox}
\section[\texorpdfstring{\ru{идите} \ru{прямо} … (idítye pryámo …)}{идите прямо … (idítye pryámo …)}]{\ru{идите} \ru{прямо}, \ru{поверните} \ru{направо} — following the answer}

\label{lesson:RU-C265-idite-pryamo}



\begin{quote}
Ask the way to the museum. (\emph{Izvinítye, kak proytí k muzéyu?}) Then say \textquotedblleft{}straight on\textquotedblright{}, \textquotedblleft{}to the right\textquotedblright{}, \textquotedblleft{}to the left\textquotedblright{}. (\emph{Pryámo. Naprávo. Nalévo.})
\end{quote}

\begin{grammarlens}[title={Grammar Lens: \ru{идите} and \ru{поверните}}]
The answer comes as instructions, with the same \textbf{-\ru{ите}} as \emph{\ru{извините}} and \emph{\ru{повторите}}:

\begin{quote}
\textbf{\ru{Идите} \ru{прямо}.} — \emph{idítye pryámo} — Go straight on. \textbf{\ru{Поверните} \ru{направо}.} — \emph{povernítye naprávo} — Turn right.
\end{quote}

\textbf{\ru{Идите}} is \emph{\ru{идти}} as an instruction. \textbf{\ru{Поверните}} comes from \emph{\ru{повернуть}}, the perfective partner of \emph{\ru{поворачивать}}: one turn, at one corner. To a friend, the \emph{-\ru{ите}} drops to \textbf{\ru{иди}}, \textbf{\ru{поверни}}.
\end{grammarlens}

\begin{grammarlens}[title={Grammar Lens: as far as, and at}]
Two small words carry a route. \textbf{\ru{до}} is \textquotedblleft{}as far as\textquotedblright{}, and the noun after it changes as in \emph{\ru{до} \ru{свидания}}: \emph{\ru{светофор}}, \textbf{\ru{до} \ru{светофора}}. \textbf{\ru{на}} is \textquotedblleft{}at\textquotedblright{} a crossroads: \emph{\ru{перекрёсток}}, \textbf{\ru{на} \ru{перекрёстке}}.

\begin{quote}
\textbf{\ru{Идите} \ru{прямо} \ru{до} \ru{светофора}, \ru{потом} \ru{поверните} \ru{налево}.} — \emph{idítye pryámo do svetofóra, potóm povernítye nalévo} — Go straight on to the lights, then turn left. \textbf{\ru{На} \ru{перекрёстке} \ru{поверните} \ru{направо}.} — \emph{na perekryóstke povernítye naprávo} — At the crossroads, turn right.
\end{quote}
\end{grammarlens}

\subsection*{The exchange}
\begin{quote}
\textbf{— \ru{Как} \ru{пройти} \ru{к} \ru{музею}? — \ru{Идите} \ru{прямо}, \ru{потом} \ru{направо}. \ru{Музей} \ru{справа}.}
\end{quote}

\emph{Muzéy správa}: \textquotedblleft{}the museum is on the right\textquotedblright{}. \textbf{\ru{Справа}} is where a thing stands; \textbf{\ru{направо}} is where you go.

\begin{cousinweb}
\textbf{\ru{направо}} is \emph{\ru{на}}, \textquotedblleft{}onto\textquotedblright{}, with the root of \emph{\ru{правый}}, \textquotedblleft{}right\textquotedblright{}: onto the right hand. EtymDB 2.1 traces \emph{\ru{правый}} to Proto-Slavic \emph{*pravŭ}, \textquotedblleft{}right, correct\textquotedblright{}, so the right hand and being right share one root. \textbf{\ru{поверните}} is on the root of \emph{\ru{вертеть}}, \textquotedblleft{}to turn\textquotedblright{}, which EtymDB takes back through Proto-Slavic \emph{*vĭrtěti} to Proto-Indo-European \emph{*wert-}: the root of Latin \emph{vertere} and English \emph{-ward}.
\end{cousinweb}

\subsection*{Your turn}
\begin{itemize}
\raggedright
  \item \emph{Say it:} \emph{Idítye pryámo do svetofóra, potóm povernítye nalévo.}
  \item \emph{Say it:} \emph{Na perekryóstke povernítye naprávo.}
  \item \emph{Say it:} the same to a friend: \emph{Idí pryámo, potóm poverní naprávo.}
  \item \emph{Answer:} \emph{Kak proytí k metró?} Give two steps, and end with \emph{správa} or \emph{sléva}
\end{itemize}

\begin{tcolorbox}[breakable,colback=teal!4,colframe=teal!35!black,title={Before you move on}]
\textquotedblleft{}Turn left\textquotedblright{}, politely? (\emph{\textbf{\ru{Поверните} \ru{налево}.}}) \textquotedblleft{}As far as the lights\textquotedblright{}? (\emph{\textbf{\ru{До} \ru{светофора}.}}) Where you go, or where it stands: \emph{\ru{направо}} or \emph{\ru{справа}}? (\emph{\textbf{\ru{Направо}}} is where you go.)
\end{tcolorbox}
\section[(review)]{Review — asking the way, and giving it}

\label{lesson:RU-C265-review}



\begin{quote}
Ask the way to the park, then to the theatre. (\emph{Kak proytí k párku? Kak proytí k teátru?})
\end{quote}

\begin{grammarlens}[title={What you've built}]
Two lessons, and a stranger's street is no longer a wall. The question: \textbf{\ru{Извините}, \ru{как} \ru{пройти} \ru{к} …?}, with the place in its \textquotedblleft{}to\textquotedblright{} form. The answer: \textbf{\ru{идите} \ru{прямо}}, \textbf{\ru{поверните} \ru{направо}} or \textbf{\ru{налево}}, \textbf{\ru{до}} something, \textbf{\ru{на}} a corner.
\end{grammarlens}

\subsection*{The exchange}
\begin{quote}
\textbf{— \ru{Извините}, \ru{как} \ru{пройти} \ru{к} \ru{театру}?} \textbf{— \ru{Это} \ru{далеко}. \ru{Вам} \ru{лучше} \ru{взять} \ru{такси}.}
\end{quote}

\emph{Vam lúchshe vzyat taksí} — \textquotedblleft{}You'd better take a taxi.\textquotedblright{} Advice from the last chapter, now with the polite \textbf{\ru{вам}}.

\subsection*{Your turn}
\begin{itemize}
\raggedright
  \item \emph{Say it:} \emph{Izvinítye, kak proytí k párku?}
  \item \emph{Say it:} \emph{Idítye pryámo do perekryóstka, potóm povernítye naprávo.}
  \item \emph{Say it:} \emph{Éto dalekó. Vam lúchshe vzyat taksí.}
  \item \emph{Answer:} a stranger asks \emph{Kak proytí k vokzálu?} Give the way from where you live, in two or three steps
\end{itemize}

\begin{tcolorbox}[breakable,colback=teal!4,colframe=teal!35!black,title={Before you move on}]
\textquotedblleft{}To the park\textquotedblright{}? (\emph{\textbf{\ru{К} \ru{парку}.}}) \textquotedblleft{}Go straight on\textquotedblright{}? (\emph{\textbf{\ru{Идите} \ru{прямо}.}}) And to a friend? (\emph{\textbf{\ru{Иди} \ru{прямо}.}})
\end{tcolorbox}
